\chapter{Belief propagation}
\label{ch:bp}

Chapter~\ref{ch:graphical} showed the posterior is a tree. This chapter turns that into an
algorithm and proves it exact.

\section{The problem}

We want $P(a_i\mid x,t)$ for each $i$, from
\begin{equation}
P(a\mid x,t) \;\propto\; \mu(a_1)\prod_{i=2}^{\nsites}\kernel(a_i\mid a_{i-1})
\prod_{i=1}^{\nsites}\like(x_i\mid a_i).
\end{equation}
By definition this means integrating out the other $\nsites-1$ variables. Discretising each
variable to $\ngrid$ values and summing naively costs $\ngrid^{\nsites-1}$ operations --- at
$\ngrid=401$, $\nsites=32$ that is about $10^{80}$, so the naive route is not merely slow but
impossible.

\section{The idea: factor out what does not depend on the summation variable}

Consider a small case, $\nsites=3$, and write $\psi_i(a_i) := \like(x_i\mid a_i)$. We want
$P(a_3\mid x)$, so we integrate over $a_1$ and $a_2$:
\[
P(a_3\mid x) \propto \psi_3(a_3)\int\!\mathrm{d}a_2 \int\!\mathrm{d}a_1\;
\mu(a_1)\psi_1(a_1)\,\kernel(a_2\mid a_1)\,\psi_2(a_2)\,\kernel(a_3\mid a_2).
\]
The inner integral over $a_1$ involves only the first three factors --- $\kernel(a_3\mid a_2)$
does not contain $a_1$, so it can be pulled out. Define
\[
\lmsg_2(a_2) := \int\!\mathrm{d}a_1\; \mu(a_1)\,\psi_1(a_1)\,\kernel(a_2\mid a_1),
\]
a function of $a_2$ alone. Then
\[
P(a_3\mid x) \propto \psi_3(a_3)\int\!\mathrm{d}a_2\;\lmsg_2(a_2)\,\psi_2(a_2)\,\kernel(a_3\mid a_2)
=: \psi_3(a_3)\,\lmsg_3(a_3).
\]

\begin{intuition}
Two integrals became two \emph{nested} integrals, each over one variable. The saving is the
whole algorithm: instead of a single sum over a grid of size $\ngrid^{2}$, we do two sums of
size $\ngrid$ each, reusing the first result. That is dynamic programming, and it works
because the factorisation lets the summand split.
\end{intuition}

\section{The general recursion}

Generalising, define the \emph{forward} (left) and \emph{backward} (right) messages
\begin{empheq}[box=\fbox]{align}
\label{eq:fwd}
\lmsg_{i}(a_i) &\;\propto\; \int\!\mathrm{d}a_{i-1}\;
\kernel(a_i\mid a_{i-1})\,\psi_{i-1}(a_{i-1})\,\lmsg_{i-1}(a_{i-1}),
&& \lmsg_1 := \mu ,\\[2pt]
\label{eq:bwd}
\rmsg_{i}(a_i) &\;\propto\; \int\!\mathrm{d}a_{i+1}\;
\kernel(a_{i+1}\mid a_i)\,\psi_{i+1}(a_{i+1})\,\rmsg_{i+1}(a_{i+1}),
&& \rmsg_{\nsites} :\equiv 1 .
\end{empheq}

The base cases matter. $\lmsg_1=\mu$ is the only place the initial law enters; without it the
prior on the first site would be silently dropped. $\rmsg_{\nsites}\equiv1$ says there is no
evidence to the right of the last site --- an improper flat function, which is fine because it
is always multiplied by something integrable.

\begin{proposition}[Exactness]
\label{prop:bp-exact}
With \eqref{eq:fwd}--\eqref{eq:bwd},
\begin{equation}
\label{eq:marginal}
P(a_i \mid x,t) \;\propto\; \lmsg_i(a_i)\,\psi_i(a_i)\,\rmsg_i(a_i),
\end{equation}
and the pairwise marginal on the edge $(i-1,i)$ is
\begin{equation}
\label{eq:pairwise}
P(a_{i-1},a_i \mid x,t) \;\propto\;
\lmsg_{i-1}(a_{i-1})\,\psi_{i-1}(a_{i-1})\;\kernel(a_i\mid a_{i-1})\;
\psi_i(a_i)\,\rmsg_i(a_i).
\end{equation}
\end{proposition}

\begin{proof}
By induction. Unrolling \eqref{eq:fwd},
\[
\lmsg_i(a_i) \propto \int\!\mathrm{d}a_{1:i-1}\;
\mu(a_1)\prod_{j=2}^{i}\kernel(a_j\mid a_{j-1})\prod_{j=1}^{i-1}\psi_j(a_j),
\]
which is exactly the joint density of everything at or left of site $i$, with the left
variables integrated out and $\psi_i$ \emph{not} yet included. Symmetrically $\rmsg_i(a_i)$ is
the joint over everything strictly right of $i$, integrated out, given $a_i$. By
Proposition~\ref{prop:separator} these two are conditionally independent given $a_i$, so their
product times the local evidence $\psi_i(a_i)$ is proportional to the marginal. The pairwise
statement is the same argument with the separator taken to be the edge rather than a node.
\end{proof}

\begin{pitfall}
Note where $\psi_i$ appears. In \eqref{eq:fwd} the local evidence at site $i-1$ is absorbed
\emph{before} the transition, so $\lmsg_i$ carries $\psi_1,\dots,\psi_{i-1}$ and $\rmsg_i$
carries $\psi_{i+1},\dots,\psi_{\nsites}$. The marginal \eqref{eq:marginal} therefore contains
each $\psi_j$ exactly once. Absorbing the local evidence \emph{after} the transition instead
would put $\psi_i$ into both messages and square it --- the belief would be sharper than the
truth, and every posterior variance too small. Getting this wrong produces plausible-looking
output, which is why the convention is pinned by a test rather than left to a comment.
\end{pitfall}

\begin{codebox}
\coderef{src/bp\_grid.py}{grid\_bp} implements \eqref{eq:fwd}--\eqref{eq:marginal} for one
sequence, returning the beliefs and the log-evidence.
\coderef{src/bp\_grid.py}{grid\_bp\_batch} does the same for a batch, sharing the site loop so
each message update is a single $(\ngrid\times\ngrid)\cdot(\ngrid\times B)$ matrix product.
The initialisations are literally \texttt{L[0] = exp(log\_mu)} and \texttt{R[-1] = 1.0}.
Pairwise statistics \eqref{eq:pairwise} are accumulated in \coderef{src/em.py}{e\_step}.
\end{codebox}

\section{Cost}

Each message update integrates one variable against the kernel: on a grid of $\ngrid$ points
that is a matrix--vector product, $O(\ngrid^2)$. There are $2\nsites$ updates, so one sequence
costs
\begin{equation}
O(\nsites \ngrid^2)
\end{equation}
instead of $O(\ngrid^{\nsites})$. At $\nsites=32$, $\ngrid=401$ this is about $5\times10^6$
operations --- microseconds --- against the $10^{80}$ of the naive route.

\begin{intuition}
The exponential became linear in the sequence length. Nothing was approximated to achieve it;
the saving is entirely from reorganising the order of integration, which the tree structure
permits and a general graph does not.
\end{intuition}

\section{Reading off the answer}

From the belief $b_i(a_i) := P(a_i\mid x,t)$ we get the posterior mean by one more integral,
\begin{equation}
\pmeansite{i} = \int\!\mathrm{d}a_i\; a_i\, b_i(a_i),
\end{equation}
and then the score from \eqref{eq:score-final}. That is the complete pipeline: messages,
beliefs, mean, score.

\begin{codebox}
\coderef{src/denoiser.py}{bp\_posterior\_mean} is the single entry point used by every
experiment --- it takes anything exposing \texttt{log\_transition\_matrix}, which covers both
the true priors of \codefile{src/priors.py} and the estimated kernels of
\codefile{src/kernels.py}. That shared interface is what makes ``exact inference under the
true kernel'' and ``inference under a learned kernel'' the same code path with a different
argument, so no comparison between them can be confounded by an implementation difference.
\end{codebox}

\section{Evidence, and why it is worth carrying}

The normalising constant of \eqref{eq:marginal} is the \emph{evidence} $\Pt(x)$. Messages are
renormalised at each step to prevent underflow, and if the discarded constants are
accumulated, their product recovers $\log \Pt(x)$ exactly.

This is not bookkeeping for its own sake. The evidence is the objective that the estimation of
Chapter~\ref{ch:estimation} maximises, and having it in closed form is what makes it possible
to check a gradient against a finite difference of the thing being optimised.

\begin{codebox}
\coderef{src/bp\_grid.py}{grid\_bp} returns \texttt{log\_evidence} alongside the beliefs.
\coderef{src/exact\_scores.py}{gaussian\_log\_evidence} gives the closed form in the Gaussian
case, which is what the general routine is tested against.
\end{codebox}

\section{The relationship to forward--backward and to Kalman}

Equations \eqref{eq:fwd}--\eqref{eq:bwd} are the forward--backward algorithm of hidden Markov
models, written for a continuous state. If the latent state is discrete the integrals become
sums and this is exactly Baum's algorithm; if everything is Gaussian, messages are determined
by two numbers each and the recursion becomes the Kalman filter and RTS smoother.

\begin{intuition}
These are not three algorithms but one, specialised by what the messages are. Discrete: a
vector. Gaussian: a mean and a variance. General continuous: a function, which must be
represented somehow --- and \emph{that} representation is the only approximation in the whole
scheme. Chapter~\ref{ch:numerics} is about it.
\end{intuition}

\section{Validation}
\label{sec:bp-validation}

Two independent checks establish that the implementation computes what this chapter claims.

\paragraph{Against brute force.} For a short chain and a coarse grid, the discretised
posterior can be enumerated exhaustively --- all $\ngrid^{\nsites}$ configurations --- and
marginalised by summation, using no messages at all. This shares no code with the recursion.
Agreement is at machine precision: posterior means to $1.6\times10^{-14}$, log-evidence to
$1.8\times10^{-15}$, and the pairwise statistic to $3.9\times10^{-16}$, with the total pairwise
mass equal to the edge count exactly.

\paragraph{Against a closed form.} On a Gaussian chain the posterior mean is known
analytically \eqref{eq:gaussian-mean}, and the recursion reproduces it to $9.2\times10^{-15}$
relative at the working resolution.

\begin{codebox}
\codefile{tests/test\_em\_bp.py} contains the enumeration check;
\codefile{tests/test\_grid\_convergence.py} and \codefile{tests/test\_score\_mean\_error\_identity.py}
contain the Gaussian and identity checks. These tests are the evidence for the numbers above
and are the reason those numbers can be quoted at all.
\end{codebox}

\section*{Sources}
\addcontentsline{toc}{section}{Sources}
The recursions of this chapter are the forward--backward algorithm for hidden Markov models
\citep{baum1970}, of which \citet{rabiner1989} is the standard exposition. In the
linear-Gaussian case they specialise to the Kalman filter \citep{kalman1960} and the
Rauch--Tung--Striebel smoother \citep{rauch1965}. \citet{cappe2005} treats the continuous-state
case in the generality needed here. That U-Net architectures can be read as approximate belief
propagation on a hierarchical model is argued by \citet{mei2024unet}.
